% Extracto LITERAL de arXiv:1409.3238 (Howlett+2015, SDSS MGS), DR7VACRSD_final.tex linea 591; bajado 2026-10-02
If we assume a $\Lambda$CDM cosmology, we are able to improve our constraints by fixing $\epsilon=0.0$ yet still allowing $\alpha$ to vary. Here we find $f\sigma_{8}=0.49^{+0.15}_{-0.14}$ and  $b\sigma_{8}=1.20^{+0.15}_{-0.15}$. This is well motivated by the Planck data, where we find that, unless we have a late time dark energy model quite different from those commonly assumed, we would expect to detect no deviation from $\epsilon=0.0$. As such this measurement is presented as our quoted, fiducial results and should be used for comparison with other $f\sigma_{8}$ results under the $\Lambda$CDM framework. However, this result should not be combined with Planck data as that would result in effectively double counting the Planck constraints. Rather, from Figure~\ref{planckAP}, we can see that combining our publicly available chains with Planck data will effectively fix $\epsilon$ and recover the fiducial results. From the same figure though we would we not recommend the usage of our results where $\alpha$ is not allowed to vary. In fact, as $\alpha$ dilates the whole correlation function, not just the BAO peak, and captures the late-time cosmological dependence of the shape of the correlation even on small scales, we would recommend that $\alpha$ be allowed to vary for any measurements of the growth of structure.
